\documentclass[12pt,a4paper]{article}
\usepackage{fontawesome5}
\providecommand{\faPencilAlt}{\faIcon{pencil-alt}}
\usepackage{styles/couleurs}
\usepackage{styles/mise_en_page}
\usepackage{styles/boites}
\usepackage{styles/maths}
\usepackage{styles/figures}
\usepackage{enumitem}
\usepackage{array}
\usepackage{tabularx}
\usepackage{amsmath}
\usepackage{amssymb}
\usepackage{tikz}
\usepackage{xparse}
\usetikzlibrary{calc,patterns,arrows.meta}
\setlist[itemize]{leftmargin=1.4em,itemsep=0.15em,topsep=0.2em}
\setlist[enumerate]{leftmargin=*,itemsep=0.25em,topsep=0.25em}
\newcommand{\lignepointillee}{\noindent\dotfill\par\vspace{0.45em}}
\newcommand{\lignesreponse}[1]{%
\begin{tcolorbox}[enhanced,colback=white,colframe=NeutreC!40,boxrule=0.5pt,
arc=1.5mm,left=3mm,right=3mm,top=2mm,bottom=1mm,before skip=2pt,after skip=2pt]
\foreach \i in {1,...,#1}{\lignepointillee}
\end{tcolorbox}}
\newcommand{\titreseq}[1]{%
\begin{center}
{\sffamily\LARGE\bfseries\color{TitreC}Séquence 4 -- Calcul littéral}\\[0.2em]
{\sffamily\large\bfseries\color{SavoirC}#1}
\end{center}}
\begin{document}

\titreseq{Progression proposée -- Écrire, simplifier et réduire}

\remarque{Durée indicative : huit séances d'apprentissage de 50 à 55 minutes,
deux interrogations courtes et une évaluation finale.

La fiche d'exercices hors classe est distribuée en début de séquence ;
les exercices sont indiqués au fur et à mesure.}

\section*{Objectifs de fin de séquence}

L'élève doit pouvoir :
\begin{itemize}
  \item traduire une phrase, une figure ou un programme par une expression littérale ;
  \item expliquer ce que représente une lettre ;
  \item calculer la valeur d'une expression pour une valeur donnée ;
  \item utiliser les conventions d'écriture du calcul littéral ;
  \item reconnaître la structure d'une expression : somme ou produit ;
  \item réduire une somme algébrique simple ;
  \item développer par distributivité simple, puis réduire si nécessaire ;
  \item distinguer une vérification sur quelques valeurs d'une preuve générale ;
  \item utiliser une expression réduite pour calculer plus efficacement.
\end{itemize}

\section*{Déroulé réaliste}

\renewcommand{\arraystretch}{1.4}
\begin{tabularx}{\textwidth}{
|>{\centering\arraybackslash}p{1.1cm}
|>{\raggedright\arraybackslash}p{4cm}
|>{\raggedright\arraybackslash}X
|>{\raggedright\arraybackslash}p{3cm}|}
\hline
\textbf{Séance} & \textbf{Support} & \textbf{Enjeu didactique} & \textbf{Production}\\
\hline
1 & Act. 1 : programmes de calcul & Passer des étapes à une expression unique ; faire émerger le besoin d'une lettre. & Expression numérique puis littérale.\\
\hline
2 & Act. 2 : bandes de carrés & Donner du sens à une formule ; faire coexister plusieurs stratégies. & Formules testées.\\
\hline
3 & Act. 3 : mur de briques & Calculer un résultat général ; faire apparaître naturellement la réduction. & Formule \(6x+5\).\\
\hline
4 & Act. 4 : conventions & Stabiliser \(3x\), \(x^2\), \(1x=x\) et la substitution avec parenthèses. & Interro 1 au début.\\
\hline
5 & Act. 5 : réduire & Regrouper les termes de même nature sans mélanger \(x\), \(x^2\) et les nombres. & Réductions simples et mixtes.\\
\hline
6 & Act. 6 : rectangles & Réactiver la distributivité ; développer puis réduire. & Deux écritures équivalentes.\\
\hline
7 & Act. 7 : programmes & Distinguer conjecture, contre-exemple et preuve. & Interro 2 au début.\\
\hline
8 & Consolidation & Automatismes, erreurs, calculs de valeurs, réduction et développement. & Synthèse du cours.\\
\hline
9 & Évaluation finale & Mobiliser écriture, substitution, simplification, réduction et preuve courte. & Évaluation sur 20.\\
\hline
\end{tabularx}

\section*{Automatismes}

Rituel de cinq minutes, quatre à six questions :
\begin{itemize}
  \item séances 1 à 3 : traduire un programme et calculer une valeur ;
  \item séances 3 à 5 : simplifier une écriture et repérer les termes semblables ;
  \item séances 5 à 7 : réduire ;
  \item séances 6 à 8 : développer puis réduire ;
  \item après la séquence : une ou deux questions chaque semaine.
\end{itemize}

\section*{Principes didactiques}

\begin{itemize}
  \item La lettre répond à une question portant sur \textbf{n'importe quelle valeur}.
  \item Plusieurs formules correctes sont conservées pour donner du sens à leur égalité.
  \item Tester une valeur peut aider à comprendre ou réfuter, mais ne prouve pas une identité.
  \item La technique vient après la situation, la verbalisation et l'écriture littérale.
  \item La résolution d'équations n'est pas l'objectif de cette séquence.
\end{itemize}

\end{document}
